% GENERATED FILE. Edit canonical lessons, then run npm run generate:books.
% canonical-source-hash: fnv1a64:6aa38da7aebe6850
% canonical-lessons: KA-C103-house, KA-C103-street, KA-C103-city, KA-C103-country, KA-C103-state, KA-C103-origin, KA-C103-origin-answer

\chapter{Where I Live}
\label{ch:ka-where-i-live}
\hlchaptermodality{103}

\begin{chapteropening}
\textbf{By the end of this chapter:} \emph{I can ask and say where somebody is from, using town, city, country and state.}

Complete the last lesson of chapter 103: I can ask and say where somebody is from.
\end{chapteropening}

\section[\texorpdfstring{\kn{ಮನೆ} (mane)}{ಮನೆ (mane)}]{\kn{ಮನೆ} (mane) — house, home}

\label{lesson:KA-C103-house}



\begin{quote}
Before the new one: say the Kannada for only one, then the Kannada for for no reason.
\end{quote}

\subsection*{You'll want to know: \kn{ಮನೆ}}
\textbf{\kn{ಮನೆ}} (\emph{mane}) — \textquotedblleft{}house, home\textquotedblright{}.

\textbf{\kn{ನಮ್ಮ} \kn{ಮನೆ}} (\emph{namma mane}) — \textquotedblleft{}our home\textquotedblright{}.

\begin{cousinweb}
\textbf{\kn{ಮನೆ}} (\emph{mane}) is native Dravidian, one word with Tamil \emph{maṉai} (DEDR 4776).
\end{cousinweb}

\begin{grammarlens}[title={What you've built}]
The place every question about you starts from.
\end{grammarlens}

\subsection*{Your turn}
\begin{itemize}
\raggedright
  \item \emph{Say it:} \emph{mane}
  \item \emph{Say it:} \emph{mane}, once more
  \item \emph{Say it:} say \emph{namma mane}
  \item \emph{Recall:} say the Kannada for only one, then the Kannada for for no reason, then say \emph{mane} again
\end{itemize}

\begin{tcolorbox}[breakable,colback=teal!4,colframe=teal!35!black,title={Before you move on}]
What does \kn{ಮನೆ} mean? (\textquotedblleft{}House, home\textquotedblright{}.) Say it once more.
\end{tcolorbox}
\section[\texorpdfstring{\kn{ಬೀದಿ} (bīdi)}{ಬೀದಿ (bīdi)}]{\kn{ಬೀದಿ} (bīdi) — street}

\label{lesson:KA-C103-street}



\begin{quote}
Before the new one: say the Kannada for for no reason, then the Kannada for home.
\end{quote}

\subsection*{You'll want to know: \kn{ಬೀದಿ}}
\textbf{\kn{ಬೀದಿ}} (\emph{bīdi}) — \textquotedblleft{}street\textquotedblright{}. It is part of an \textbf{\kn{ವಿಳಾಸ}}, \textquotedblleft{}address\textquotedblright{}.

\begin{cousinweb}
\textbf{\kn{ಬೀದಿ}} (\emph{bīdi}) comes from Sanskrit \emph{vīthi}, \textquotedblleft{}road, street; line, row\textquotedblright{}, reshaped to Kannada sounds: a tadbhava (CDIAL 12050). The Sanskrit \emph{v-} became Kannada \emph{b-}.
\end{cousinweb}

\begin{grammarlens}[title={What you've built}]
One line of an address.
\end{grammarlens}

\subsection*{Your turn}
\begin{itemize}
\raggedright
  \item \emph{Say it:} \emph{bīdi}
  \item \emph{Say it:} \emph{bīdi}, once more
  \item \emph{Say it:} say \emph{bīdi}
  \item \emph{Recall:} say the Kannada for for no reason, then the Kannada for home, then say \emph{bīdi} again
\end{itemize}

\begin{tcolorbox}[breakable,colback=teal!4,colframe=teal!35!black,title={Before you move on}]
What does \kn{ಬೀದಿ} mean? (\textquotedblleft{}Street\textquotedblright{}.) Say it once more.
\end{tcolorbox}
\section[\texorpdfstring{\kn{ನಗರ} (nagara)}{ನಗರ (nagara)}]{\kn{ನಗರ} (nagara) — city}

\label{lesson:KA-C103-city}



\begin{quote}
Before the new one: say the Kannada for home, then the Kannada for street.
\end{quote}

\subsection*{You'll want to know: \kn{ನಗರ}}
\textbf{\kn{ನಗರ}} (\emph{nagara}) — \textquotedblleft{}city\textquotedblright{}. \textbf{\kn{ಊರು}} is any town, and home; \textbf{\kn{ನಗರ}} is the big city.

\begin{cousinweb}
\textbf{\kn{ನಗರ}} (\emph{nagara}) is Sanskrit \emph{nagara}, \textquotedblleft{}town\textquotedblright{}, taken over whole as a learned loan, a tatsama (CDIAL 6924).
\end{cousinweb}

\begin{grammarlens}[title={What you've built}]
Bigger than a town.
\end{grammarlens}

\subsection*{Your turn}
\begin{itemize}
\raggedright
  \item \emph{Say it:} \emph{nagara}
  \item \emph{Say it:} \emph{nagara}, once more
  \item \emph{Say it:} say \emph{nagara}
  \item \emph{Recall:} say the Kannada for home, then the Kannada for street, then say \emph{nagara} again
\end{itemize}

\begin{tcolorbox}[breakable,colback=teal!4,colframe=teal!35!black,title={Before you move on}]
What does \kn{ನಗರ} mean? (\textquotedblleft{}City\textquotedblright{}.) Say it once more.
\end{tcolorbox}
\section[\texorpdfstring{\kn{ದೇಶ} (dēśa)}{ದೇಶ (dēśa)}]{\kn{ದೇಶ} (dēśa) — country}

\label{lesson:KA-C103-country}



\begin{quote}
Before the new one: say the Kannada for street, then the Kannada for city.
\end{quote}

\subsection*{You'll want to know: \kn{ದೇಶ}}
\textbf{\kn{ದೇಶ}} (\emph{dēśa}) — \textquotedblleft{}country\textquotedblright{}.

\textbf{\kn{ನಿಮ್ಮ} \kn{ದೇಶ} \kn{ಯಾವುದು}?} (\emph{nimma dēśa yāvudu?}) — \textquotedblleft{}what is your country?\textquotedblright{}

\begin{cousinweb}
\textbf{\kn{ದೇಶ}} (\emph{dēśa}) is Sanskrit \emph{deśa}, \textquotedblleft{}point, region; country\textquotedblright{}, taken over whole as a learned loan, a tatsama (CDIAL 6547).
\end{cousinweb}

\begin{grammarlens}[title={What you've built}]
The widest home of all.
\end{grammarlens}

\subsection*{Your turn}
\begin{itemize}
\raggedright
  \item \emph{Say it:} \emph{dēśa}
  \item \emph{Say it:} \emph{dēśa}, once more
  \item \emph{Say it:} ask \emph{nimma dēśa yāvudu?}
  \item \emph{Recall:} say the Kannada for street, then the Kannada for city, then say \emph{dēśa} again
\end{itemize}

\begin{tcolorbox}[breakable,colback=teal!4,colframe=teal!35!black,title={Before you move on}]
What does \kn{ದೇಶ} mean? (\textquotedblleft{}Country\textquotedblright{}.) Say it once more.
\end{tcolorbox}
\section[\texorpdfstring{\kn{ರಾಜ್ಯ} (rājya)}{ರಾಜ್ಯ (rājya)}]{\kn{ರಾಜ್ಯ} (rājya) — state (of a country)}

\label{lesson:KA-C103-state}



\begin{quote}
Before the new one: say the Kannada for city, then the Kannada for country.
\end{quote}

\subsection*{You'll want to know: \kn{ರಾಜ್ಯ}}
\textbf{\kn{ರಾಜ್ಯ}} (\emph{rājya}) — \textquotedblleft{}state\textquotedblright{}. \textbf{\kn{ಕರ್ನಾಟಕ} \kn{ರಾಜ್ಯ}} (\emph{karnāṭaka rājya}) is the state where Kannada is spoken.

\begin{cousinweb}
\textbf{\kn{ರಾಜ್ಯ}} (\emph{rājya}) is Sanskrit \emph{rājya}, \textquotedblleft{}kingship, kingdom\textquotedblright{}, taken over whole as a learned loan, a tatsama (CDIAL 10694).
\end{cousinweb}

\begin{grammarlens}[title={What you've built}]
The state where Kannada lives.
\end{grammarlens}

\subsection*{Your turn}
\begin{itemize}
\raggedright
  \item \emph{Say it:} \emph{rājya}
  \item \emph{Say it:} \emph{rājya}, once more
  \item \emph{Say it:} say \emph{karnāṭaka rājya}
  \item \emph{Recall:} say the Kannada for city, then the Kannada for country, then say \emph{rājya} again
\end{itemize}

\begin{tcolorbox}[breakable,colback=teal!4,colframe=teal!35!black,title={Before you move on}]
What does \kn{ರಾಜ್ಯ} mean? (\textquotedblleft{}State (of a country)\textquotedblright{}.) Say it once more.
\end{tcolorbox}
\section[\texorpdfstring{\kn{ನೀವು} \kn{ಎಲ್ಲಿಂದ} … (nīvu ellinda …)}{ನೀವು ಎಲ್ಲಿಂದ … (nīvu ellinda …)}]{\kn{ನೀವು} \kn{ಎಲ್ಲಿಂದ} \kn{ಬಂದಿದ್ದೀರಿ}? (nīvu ellinda bandiddīri?) — where are you from?}

\label{lesson:KA-C103-origin}



\begin{quote}
Ask \textbf{\kn{ನಿಮ್ಮ} \kn{ಊರು} \kn{ಯಾವುದು}?} Then say city, country and state.
\end{quote}

\begin{grammarlens}[title={Grammar Lens: -\kn{ಇಂದ} means “from”}]
Change Chapter 5's \textbf{\kn{ಬೆಂಗಳೂರಿನಲ್ಲಿ}} (\emph{Beṅgaḷūrinalli}, “in Bengaluru”):

\begin{quote}
\textbf{\kn{ಬೆಂಗಳೂರಿನಿಂದ}} — \emph{Beṅgaḷūrininda} — \textbf{from Bengaluru}
\end{quote}

\textbf{-\kn{ಇಂದ}} (\emph{-inda}) marks a starting place. \textbf{\kn{ಎಲ್ಲಿ}} (\emph{elli}, “where?”) becomes \textbf{\kn{ಎಲ್ಲಿಂದ}} (\emph{ellinda}), “from where?”
\end{grammarlens}

\subsection*{You'll want to know: ask it as one respectful frame}
\begin{quote}
\textbf{\kn{ನೀವು} \kn{ಎಲ್ಲಿಂದ} \kn{ಬಂದಿದ್ದೀರಿ}?}
\end{quote}

\begin{quote}
\emph{nīvu ellinda bandiddīri?} — “Where are you from?”
\end{quote}

Keep \textbf{\kn{ಬಂದಿದ್ದೀರಿ}} (\emph{bandiddīri}) inside this fixed frame.

\subsection*{Your turn}
\begin{itemize}
\raggedright
  \item \emph{Say it:} \emph{Beṅgaḷūru} — Bengaluru; \emph{Beṅgaḷūrininda} — from Bengaluru
  \item \emph{Ask:} \textbf{\kn{ನೀವು} \kn{ಎಲ್ಲಿಂದ} \kn{ಬಂದಿದ್ದೀರಿ}?}
  \item \emph{Recall:} ask for a home town, then say city, country and state
\end{itemize}

\begin{tcolorbox}[breakable,colback=teal!4,colframe=teal!35!black,title={Before you move on}]
Ask “Where are you from?” Which ending turns \textbf{\kn{ಎಲ್ಲಿ}} into “from where”? (\textbf{-\kn{ಇಂದ}}, “from”.)
\end{tcolorbox}
\section[\texorpdfstring{\kn{ನಾನು} \kn{ಬೆಂಗಳೂರಿನಿಂದ} … (nānu …)}{ನಾನು ಬೆಂಗಳೂರಿನಿಂದ … (nānu …)}]{\kn{ನಾನು} \kn{ಬೆಂಗಳೂರಿನಿಂದ} \kn{ಬಂದಿದ್ದೇನೆ}. (nānu Beṅgaḷūrininda bandiddēne) — I am from Bengaluru}

\label{lesson:KA-C103-origin-answer}



\begin{quote}
Ask \textbf{\kn{ನೀವು} \kn{ಎಲ್ಲಿಂದ} \kn{ಬಂದಿದ್ದೀರಿ}?} Then say \textbf{\kn{ಬೆಂಗಳೂರಿನಿಂದ}}.
\end{quote}

\subsection*{You'll want to know: the answer}
\begin{quote}
\textbf{\kn{ನಾನು} \kn{ಬೆಂಗಳೂರಿನಿಂದ} \kn{ಬಂದಿದ್ದೇನೆ}.}
\end{quote}

\begin{quote}
\emph{nānu Beṅgaḷūrininda bandiddēne.} — “I am from Bengaluru.”
\end{quote}

\textbf{\kn{ಎಲ್ಲಿಂದ}} asks; \textbf{\kn{ಬೆಂಗಳೂರಿನಿಂದ}} answers. Keep the whole answer as one frame.

\subsection*{The exchange: a second shape to recognise}
Recognise this attested short exchange:

\begin{quote}
\textbf{\kn{ನೀನು} \kn{ಕರ್ನಾಟಕದವನೇ}?} — “Are you from Karnataka?”
\end{quote}

\begin{quote}
\textbf{\kn{ಹೌದು} \kn{ಸರ್}, \kn{ಬೆಂಗಳೂರಿನಿಂದ}.} — “Yes sir, from Bengaluru.”
\end{quote}

\textbf{\kn{ಕರ್ನಾಟಕದವನೇ}} describes a male addressee, so keep it for recognition. Produce the respectful \textbf{\kn{ಎಲ್ಲಿಂದ}} question and an \textbf{-\kn{ಇಂದ}} answer.

\subsection*{Your turn}
\begin{itemize}
\raggedright
  \item \emph{Ask:} \textbf{\kn{ನೀವು} \kn{ಎಲ್ಲಿಂದ} \kn{ಬಂದಿದ್ದೀರಿ}?}
  \item \emph{Answer:} \textbf{\kn{ನಾನು} \kn{ಬೆಂಗಳೂರಿನಿಂದ} \kn{ಬಂದಿದ್ದೇನೆ}.}
  \item \emph{Recognise:} \textbf{\kn{ಕರ್ನಾಟಕದವನೇ}?} asks whether a male addressee is from Karnataka
\end{itemize}

\begin{tcolorbox}[breakable,colback=teal!4,colframe=teal!35!black,title={Before you move on}]
Ask “Where are you from?” Answer “I am from Bengaluru.” Then recognise the Karnataka relation question.
\end{tcolorbox}
